% Transcription — fonds Alexandre Grothendieck, shelfmark 37, batch 5 (pages 81-99).
% Established from the facsimile archives/batches/37/batch-05.pdf (a copy of
% 37.pdf fetched through the project's relay, grothendieck-relay.onrender.com;
% PDF page k = archivists' page 80+k, no cover sheet), per the skill
% transcribe-grothendieck. The folder is 99 pages in five batches (1-20, 21-40,
% 41-60, 61-80, 81-99); this is the fifth and last, nineteen pages, done in
% parallel with the others.
%
% Pass: Opus 5.5 (claude-opus-5-5), 2026-09-26 — first pass, unchecked against the pages by a human.
%
% Pages skipped: none outright. Pages 83 and 98 are tan covers inscribed top
% right, « Réédition EGA I » (ink, underlined, a scribbled-out word below it)
% and « EGA II » (pencil): no \page marker, their words become section
% headings with a \note (hand.md §5; the same treatment as the SGA covers of
% batches 1, 3 and 4).
% Pages noted only (not re-set): 84-91 (typescript of EGA I), 95 (Dieudonné).
% Pages transcribed: 81, 82, 92, 93, 94, 96, 97, 99.
%
% Decisions, item by item, with the evidence:
% - Pages 81-82: his own manuscript notes (ink; p. 81 on IHES letterhead),
%   transcribed in full. They continue the SGA 7 run of batch 4 (cover
%   « SGA 7 », p. 72; pp. 76-80): p. 81 items for XXI (appendix « de Katz »,
%   torsion in crystalline cohomology), XXII (a reference to Deligne's proof
%   of Katz's congruence theorem) and VI or XIII (ordinary quadratic
%   singularities, openness of the locus where a flat morphism of finite
%   presentation is smooth or has such a singularity); p. 82, headed « SGA 7
%   XVIII », page/line places with the references to put there (SGA 5 IV,
%   Exp XIV, XVI, XII) — the same form as p. 80 (batch 4), which also lists
%   places in XVIII. The section heading is therefore « SGA 7 (suite) ».
% - Pages 84-88, 90: a typescript « TABLE DES MATIERES », typed pagination
%   « - 538 - » to « - 543 - » (the leaves lie in the order 538, 541, 539,
%   540, 542, 543; « 542 » is typed over another figure), running from the
%   Avant-propos, Introduction, Chapitre 0 (§§ 1-7) and Chapitre I (§§ 1-10)
%   to an « APPENDICE .- Ultraschémas et espaces algébriques de Serre », a
%   bibliography and two indexes. DECISION: typescript of a PUBLISHED text,
%   one French \note per page, not re-set. Evidence: it is the table of the
%   new edition of EGA I (Grundlehren 166, Springer 1971) — the new
%   Chapitre 0 §§ 1-7 with 2.7 Quasi-homéomorphismes, 2.8 Espaces de
%   Jacobson, 2.9 Espace sobre associé, and Chapitre I with the new § 7
%   « Ensembles constructibles dans les schémas » and the Serre appendix —
%   whose structure it matches as far as it is known to this pass (not
%   collated line by line with the printed book). His marks: on p. 84 a pencil
%   wavy stroke cancelling the two superseded lines « 2.4. Fonctions
%   constructibles » (also struck through at the machine) and « 2.5. », with
%   a small pencil bracket in the left margin; on p. 90 a pencil box
%   « Morphismes affines » with a line leading to below 10.15, i.e. a new
%   subsection to be added after 10.15. Whose pencil is not established (it is
%   given in full, presumably his, not proven). Pages 85-88 carry nothing
%   handwritten (typed x-outs only).
% - Pages 89, 91: two unpaginated typed leaves (« -   - ») of the same
%   volume's « Index terminologique », letters S to V (p. 91 « S-schéma
%   localement quasi-fini ... Spectre maximal ... », p. 89 « Spectre formel
%   ... Valeur »), with references of the form « 0,3.1.6 », « I,10.4.2 »,
%   « I,App.1.5 ». DECISION: part of the published EGA I (1971) apparatus,
%   one \note each, not re-set; nothing handwritten on them.
% - Pages 92-93: his typed letter « Cher Dieudonné », dated by hand
%   « 27.8.1967 » (top right of p. 92), on EGA I 10.10.5 (a counterexample
%   without noetherian hypotheses) and, on p. 93, a « canular » in EGA IV
%   17.1.6 (i) and the critiques of Samuel and Illusie on « mon projet de
%   par. 0 pour la réédition de EGA I ». DECISION: his letter, transcribed in
%   full. Evidence: Dieudonné's reply (p. 95) thanks « de ta lettre du 27/8 »
%   and answers exactly these points (the counterexample, 10.10.5 a) and c),
%   « tes remarques sur l'Introduction »); p. 93 is the verso of the same
%   sheet (the date of p. 92 shows through mirrored at its top left) and
%   carries the continuation; it breaks off mid-sentence (« le sens du mot »),
%   the rest is not in the folder's batch. The machine's missing symbols
%   (brackets of B[[T]], the element and inequality signs, Omega of
%   Omega^1_{X/Y}) are inked on p. 92 and left blank on p. 93, where the
%   Omega is \supplied{}. Two long oblique strokes cross the last paragraph
%   of p. 93; whose is not established.
% - Page 94: a single line in his hand on IHES letterhead, « 10.12.
%   Stabilité par lim (ou non ?) » (lim with a leftward arrow under it),
%   transcribed; 10.12 is « Morphismes adiques de schémas formels » in the
%   table of p. 90.
% - Page 95: a typed letter of J. Dieudonné, signed « J A Dieudonné » in blue
%   ink, « Nice , 11 Septembre 1967 », « Mon cher Grothendieck »: EGA I
%   10.10.5 (still not convinced that a) implies c) for modules of finite
%   presentation; the implication b) => a) used in the construction of G is
%   false; « canulars Artin-Rees »), the Introduction to be written after
%   § 10, and practical matters (Serre's paper to go to the Presses through
%   Mlle Rolland, Segal's memoirs). Another person's letter: one French
%   \note, not re-set. His marks, given in full: black ink in the left margin
%   « cf. EGA IV 18.3.2.1. » (his, with good likelihood: his reply of p. 97
%   cites exactly « EGA IV 18.3.2.1 »), and a double vertical stroke in the
%   margin against the last paragraph (hand not established). The yellow and
%   green underlinings of O_X, G, F and the blue arrows are the writer's
%   typographical marks, not his.
% - Pages 96-97: his typed letter « Cher Dieudonné », dated by hand
%   « 15.9.67. », answering « ta lettre du 11 Septembre »: a Lemme on the
%   number of generators of a module that is locally finitely generated,
%   uniform p and q in 10.10.5, the corollary a)-c) without noetherian
%   hypotheses, EGA IV 18.3.2.1 to become a corollary after 0_I 7.2.9, and
%   the fate of the modifications proposed for I 10.11. DECISION: his letter,
%   transcribed in full; p. 97 continues p. 96 (« ... par mp éléments / où m
%   est le nombre des f_i »), no signature on the leaf. The machine's blanks
%   for arrows, the element and inequality signs and the tilde were never
%   inked: they are restored in the formulas and said once in a \note;
%   one blank the text does not settle (« m _ m_i ») is left and noted.
% - Page 99: his manuscript notes (ink, IHES letterhead), items (a)-(d) for
%   the re-edition after the « EGA II » cover: EGA IV 4.8 to be transferred,
%   an appendix on formal schemes, EGA II 6.5.10, 6.5.11, 6.5.7 (with a
%   reference to EGA IV 21), EGA IV 19.4.11 to be brought up « revu et
%   augmenté », 7.4.2. Transcribed in full.
%
% His pagination and numbered runs: none of his own on his manuscript leaves.
% His lists are keyed to the pages and lines of a typescript of SGA 7 XVIII
% (p. 82: pp. 10-68) and to numbers of EGA (p. 99: items (a)-(d), the
% letters ringed). The EGA I typescript carries typed pagination 538-543 (pp.
% 84-88, 90); its index leaves (pp. 89, 91) are unpaginated. Dates on the
% leaves: « 27.8.1967 » (p. 92, his hand); « Nice , 11 Septembre 1967 »
% (Dieudonné, p. 95, which cites « ta lettre du 27/8 »); « 15.9.67. » (p. 96,
% his hand, which cites « ta lettre du 11 Septembre »).
% ALIGNMENT: batches 1, 3 and 4 existed when this file was finished and were
% followed (cover treatment; \note format for published typescripts and
% other people's letters; his typed letters transcribed in full; roman exposé
% numbers plain and « l.~$-n$ » for line references, as in batch 4). Batch 2
% did not exist.
%
% What the batch is about. The end of the SGA 7 run (appendices to XXI and
% XXII on crystalline torsion and Katz's congruence formula; ordinary
% quadratic singularities; reference corrections in XVIII), then the
% re-edition of EGA I and II: the typescript table of contents and part of
% the index of the new EGA I, and the Grothendieck-Dieudonné exchange of
% August-September 1967 on EGA I 10.10.5 without noetherian hypotheses — a
% counterexample (A = B[[T]], M_n = A_n/J_n with J_n = (X_0, X_1 T, ...,
% X_n T^n)), a lemma bounding generators of a locally finitely generated
% module, the uniformity of presentations in n, and the corollary that F is
% locally free of finite type iff it is a limit of compatible locally free
% F_n iff it is M~ with M projective of finite type — together with a
% « canular » in EGA IV 17.1.6 (i) (is « formellement lisse » local? reduced
% to: is a module locally projective on Spec(A) projective?). Last, a list
% of items for EGA II.
%
% Notation fixed once: roman exposé numbers plain; \underline{O} for the
% typescript's underlined O; \widetilde{M} for the tilde left blank; \lim
% for the typed « lim » (p. 96), \varprojlim for his arrowed lim (p. 94).
% Machine overstrikes (xxx) are not recorded.
% Letters lost where the typed lines run off the right edge of the leaf
% (pp. 92, 93, 96, 97) are \supplied{}.
%
% Hardest pages: 81 (fast ink; « déformant », « donne », the word before
% « ouvert ») and 99 (several words not read). Apparatus: 9 \ill{},
% 16 \uncertain{} (counted on the finished file).
%
% Readings to check first: on p. 81 « qu'en déformant », « donne », « car.
% rés. 2 », « l'ens. » and the word before « ouvert »; on p. 82 « Exp », the
% struck word of the first line and « 15 ? »; on p. 99 « écrit », the
% third line of (a), the word after « pour » in « Appendice pour ... »,
% « dans le ... », « gén. décomposé », the parenthesis « (... !) » and the
% two words after « Admettre ».

\documentclass[11pt,a4paper]{article}
\input{../preamble/grothendieck}

\folder{37}
\batch{5}
\pages{81}{99}
\dating{1967-1971}
\watermark{Édition de démonstration}
\foldertitle{Rééditions SGA, EGA [SGA 4 à SGA 7, EGA I et EGA II] : tapuscrits et copies de tapuscrit annoté (s.d.), lettres (1967-1971, s.d.), notes manuscrites (s.d.)}

\begin{document}

\section*{SGA 7 (suite)}

\note{les pages 81 et 82 continuent les notes pour SGA 7 qui suivent le
feuillet de garde « SGA 7 » (page 72, lot 4) ; le titre est repris de ce
feuillet}

\page{81}

\note{à l'encre, sur une feuille à en-tête de l'IHÉS ; les numéros d'exposé
sont soulignés par lui}

\begin{itemize}
\item Appendice à XXI de Katz : torsion en coh. cristalline
\item[XXII] Référence à la démonstration de \emph{Deligne} du th. de
congruence de Katz (généralisé)
\item[VI \emph{ou} XIII] Vérifier qu'on y trouve le fait qu'en
\uncertain{déformant} une singularité quadratique ordinaire, ---
\uncertain{donne} au plus autant (cf. ci-dessus cas de \uncertain{car. rés.} 2
et d'une relation paire). En d'autres termes, \uncertain{l'ens.} des pts où un
morphisme plat \struck{de} \add{loc.} prés. finie est lisse ou a une
singularité quadratique ordinaire est \ill{} \emph{ouvert}.
\end{itemize}

\note{le « VI » est récrit d'un trait épais ; « ou XIII » est ajouté dessous,
dans la marge. « loc. » est écrit au-dessus de « de », biffé}

\page{82}

\note{à l'encre, en haut d'une feuille blanche ; rien d'autre sur la page}

SGA 7 XVIII \quad p.~10 l.~2 \quad réf à \struck{\ill{}} \struck{Exp} SGA 5 IV

\begin{itemize}
\item p.~17 \quad l.~9 \quad réf \uncertain{Exp} XIV\note{le 1 de 17 est
récrit}
\item p.~25 \quad l.~$-2$ \quad --- SGA 5 IV
\item p.~34 \quad l.~4 \quad --- XVI
\item p.~44 \quad l.~1
\item p.~64 \quad l.~$-7$, l.~$-2$ \quad --- XVI
\item p.~66 \quad l.~6 \quad --- XII
\item l.~$-5$ \quad \uncertain{15} ?\note{lu aussi
« 18 » ; sans trait au-dessus ni au-dessous, à la différence des numéros
d'exposé voisins}
\item p.~68 \quad l.~4 \quad XVI
\item l.~$-6$ \quad XII
\end{itemize}

\section*{Réédition EGA I}

\note{inscrit à l'encre et souligné en haut à droite du feuillet de garde
(page 83), qui ne porte rien d'autre qu'un mot raturé au-dessous,
\struck{\ill{}} ; vraisemblablement de sa main. Le feuillet ne reçoit pas de
numéro de page}

\page{84}

\note{feuillet dactylographié, paginé « - 538 - » à la machine : première page
de la « TABLE DES MATIERES » de la nouvelle édition d'EGA I (Springer,
Grundlehren 166, 1971), texte publié, non reproduit : Avant-propos,
Introduction, Chapitre 0 « Préliminaires », § 1 « Foncteurs représentables »
(1.1 à 1.7), § 2 « Compléments de topologie » (2.1 à 2.10), § 3 « Compléments
sur les faisceaux » (3.1 à 3.7). L'identification repose sur la structure de
la table (nouveau Chapitre 0 avec quasi-homéomorphismes, espaces de Jacobson,
espace sobre associé), non sur une collation ligne à ligne. Marques
manuscrites, au crayon, de main non établie (vraisemblablement la sienne) :
un trait ondulé biffe les deux lignes périmées « 2.4. Fonctions
constructibles » (déjà barrée à la machine) et « 2.5. », qui précèdent les
lignes retapées « 2.4. Ensembles constructibles dans les espaces
noethériens » et « 2.5. Fonctions constructibles » ; en marge de gauche, un
petit crochet en regard de ces lignes}

\page{85}

\note{même table, page « - 541 - » (les feuillets sont rangés dans l'ordre
538, 541, 539, 540), non reproduite : Chapitre I, fin du § 3 (3.8 à 3.10,
morphismes universellement ouverts, générisants, submersifs), § 4
« Sous-schémas et morphismes d'immersion. Schémas réduits » (4.1 à 4.6), § 5
« Morphismes séparés. Critères valuatifs » (5.1 à 5.5), § 6 « Conditions de
finitude relatives » (6.1 à 6.11). Rien de manuscrit}

\page{86}

\note{même table, page « - 539 - », non reproduite : Chapitre 0, fin du § 3
(3.8, 3.9), § 4 « Espaces annelés » (4.1 à 4.5), § 5 « Faisceaux
quasi-cohérents et faisceaux cohérents » (5.1 à 5.7), § 6 « Compléments
d'algèbre commutative » (6.1 à 6.8), § 7 « Compléments d'algèbre
topologique » (7.1 à 7.6). Rien de manuscrit}

\page{87}

\note{même table, page « - 540 - », non reproduite : Chapitre 0, fin du § 7
(7.7, 7.8) ; Chapitre I « Le langage des schémas », Sommaire, § 1 « Schémas
affines » (1.1 à 1.7), § 2 « Schémas et morphismes de schémas » (2.1 à 2.8),
§ 3 « Produit et somme de schémas. Changement de base » (3.1 à 3.7). Rien de
manuscrit}

\page{88}

\note{même table, page « - 542 - » (le numéro frappé sur un autre chiffre),
non reproduite : Chapitre I, § 7 « Ensembles constructibles dans les
schémas » (7.1 à 7.3), § 8 « Applications rationnelles » (8.1 à 8.5), § 9
« Foncteurs représentables élémentaires dans la théorie des schémas » (9.1
à 9.10), § 10 « Schémas formels » (10.1 à 10.8). Rien de manuscrit}

\page{89}

\note{feuillet dactylographié non paginé (« -   - ») de l'« Index
terminologique » du même volume, de « Spectre formel d'un anneau admissible :
I,10.1.2 » à « Valeur en un point d'une section d'un $O_X$-Module sur un
espace localement annelé : 0,4.1.9 » ; il fait suite à la page 91. Texte
publié, non reproduit ; rien de manuscrit}

\page{90}

\note{même table des matières, page « - 543 - », la dernière, non reproduite :
Chapitre I, § 10, 10.9 à 10.15 (10.12 « Morphismes adiques de schémas
formels », 10.13 « Morphismes de type fini », 10.15 « Schémas formels
séparés »), « APPENDICE .- Ultraschémas et espaces algébriques de Serre »
(App. 1, App. 2), Bibliographie, Index des notations, Index terminologique.
Marque manuscrite, au crayon, de main non établie (vraisemblablement la
sienne) : dans la marge de droite, en regard de 10.12 et 10.13, un cadre
contenant « Morphismes affines », d'où part un trait qui aboutit sous la
ligne 10.15 : un numéro à ajouter après 10.15}

\page{91}

\note{feuillet dactylographié non paginé du même « Index terminologique », de
« S-schéma localement quasi-fini sur S : I,6.11.3 » à « Spectre maximal d'un
anneau de Jacobson : I,App.1.3 » ; la page 89 en est la suite. Texte publié,
non reproduit ; rien de manuscrit}

\page{92}

\note{lettre dactylographiée de lui, datée à la main en haut à droite ;
les crochets de $B[[T]]$ et de $k[(X_i)_{i\geq 0}]$, les signes $\in$,
$\leq$, $\geq$ manquaient à la machine et sont complétés à l'encre. Les
surfrappes de la machine ne sont pas relevées}

27.8.1967

Cher Dieudonné,

Tu as bien fait de tiquer sur la généralisation que je suggérai de EGA I
10.10.5. Je pense que l'analogue nonnoethérien de l'équivalence de a) et c)
doit sortir sans mal, mais en tous cas ces conditions ne son\supplied{t} pas impliquées
par b). Pour le voir, soit $A = B[[T]]$, et définissons
\[
M_n = A_n/J_n ,
\]
où $A_n = B[T]/(T^{n+1})$, $J_n = (X_0, X_1 T, \ldots, X_n T^n)$,
en choisissant l'anneau $B$ et les $X_n \in B$ de telle façon que les
générateurs écrits de $J_n$ soient en nombre minimal. IL suffit pour ceci de
prendre pour $B$ un anneau de polynômes
\[
B = k[(X_i)_{i \geq 0}] ,
\]
on vérifie en effet que les générateurs écrits donnent alors dans $J_n/TJ_n$
des éléments linéairement indépendants sur $A_0 = A/TA = B$. En effet, si les
$a_i \in B$ $(0 \leq i \leq n)$ sont tels que
\[
a_0 X_0 + \cdots + a_n X_n T^n = F_0 X_0 T + \cdots + F_n X_n T^{n+1} ,
\]
où les $F_n$ sont dans $B[T] = A[T,(X_i)]$, on voit pour tout $i$, $0 \leq i \leq n$,
divisant par l'idéal engendré par les $X_j$ avec $j \neq i$, qu'on a une
relatio\supplied{n}
\[
a_i X_i T^i = F'_i X_i T^{i+1}
\]
dans $k[X_i, T]$,
où $F'_i$ est l'image de $F_i$. Comme $X_i$ et $T$ sont des éléments
réguliers de cet anneau, on en conclut que
\[
a_i = F'_i T
\]
dans $k[X_i, T]$,
d'où évidemment $a_i = 0$, cqfd.
\note{la parenthèse fermante de $J_n$ est coupée au bord du feuillet ;
« $B[T] = A[T,(X_i)]$ » est ainsi à la machine, le contexte attendrait
$k[T,(X_i)]$}

Je suggère donc de rédiger la nouvelle proposition 10.10.5 en disant qu'on a
deux conditions équivalentes impliquant une troisième condition, et que cette
dernière équivaut aux deux autres dans le cas noethérien. Il serait peut-être
utile de donner aussi le contre-exemple qui précède.

\page{93}

\note{suite de la même lettre, au verso du même feuillet ; les blancs laissés
par la machine n'y sont pas complétés}

J'ai découvert un canular ennuyeux dans EGA IV 17.1.6 (i), concernan\supplied{t} le cas
« formellement lisse ». En effet, le démonstration via (16.5.17) ne marche que
si on suppose \supplied{$\Omega$}$^1_{X/Y}$ de présentation finie, par exemple
si $f$ est localement de type fini ! J'ignore donc si la notion « formellement
lisse » est vraiment locale en haut ! On peut montrer que la question équivaut
à la suivante : si $A$ est un anneau commutatif, et $M$ un $A$-Module (pas
nécessairement de présentation finie) qui est « localement projectif sur
$\operatorname{Spec}(A)$ », i.e. tel qu'il existe des $f_i$ \supplied{$\in$}
$A$ engendrant l'idéal unité, avec $M_{f_i}$ projectif sur $A_{f_i}$, alors
$M$ est-il un $A$-Module projectif ? (Je me rappelle avoir posé la question à
D. Lazard, mais je ne me rappelle pas s'il l'avait résolue dans un sens ou
dans l'autre. Je crois que non). Il faudrait donc à la prochaine occasion
faire un errata pour 17.1.6, dans lequel on signalerait également la
formulation équivalente que je viens de signaler pour la question du
caractère local de la lissité formelle. --- Si la réponse devait être
négative, il y aurait lieu de considérer aussi la notion de morphisme
localement formellement lisse. En fait, dans tous les cas que j'ai
rencontrés, on prouve toujours directement la lissité formelle globale, telle
qu'elle est définie dans le texte.

J'ai fait lire mon projet de par.~0 pour la réédition de EGA I à Samuel et
Illusie, dont je te transmets certaines critiques. Page 4, au lieu de
« quasi-totalité », il est plus prudent de dire « plupart ». Page 8,
j'utilise tantôt la terminologie « idéal racine », tantôt « idéal radical »
et j'écris $\operatorname{rad}(J)$, --- il faut unifier. Page 9, ligne 4,
style vicieux, lire « comme anneaux de valeurs pour les coordonnées de
solutions \ldots ». Page 10, donner en note de bas de page le sens du mot
\note{la lettre s'interrompt ici, au bas du verso ; la suite n'est pas dans
ce lot. Deux longs traits obliques, de main non établie, traversent tout ce
dernier alinéa}

\page{94}

\note{au crayon, en haut d'une feuille à en-tête de l'IHÉS ; rien d'autre sur
la page. 10.12 est, dans la table de la page 90, « Morphismes adiques de
schémas formels »}

10.12. \quad Stabilité par $\varprojlim$ (ou non ?)

\page{95}

\note{lettre dactylographiée de J. Dieudonné, signée à l'encre bleue
« J A Dieudonné », « Nice , 11 Septembre 1967 », « Mon cher Grothendieck »,
non reproduite. Il remercie de la lettre du 27/8 (pages 92-93) et du
contre-exemple, mais n'est « toujours pas convaincu » que, pour les Modules
de présentation finie sans condition noethérienne, la condition a) de
(10.10.5) entraîne c) : la construction de $G$ dans le texte actuel utilise
l'implication b) $\Rightarrow$ a), fausse dans le cadre de la présentation
finie, et l'on retombe sur « les canulars Artin-Rees » faute de suites
exactes $A_n^p \to A_n^q \to M_n \to 0$ avec $p$ et $q$ fixes ; « j'ai des
doutes sur la validité du théorème ». Il rédigera l'Introduction quand il
aura terminé le § 10 ; il accepte de prendre le papier de Serre (à remettre à
Mlle Rolland pour les Presses) et les mémoires de Segal. Les soulignements
jaunes et verts et les flèches à l'encre bleue sont du scripteur. De sa main,
à l'encre noire, dans la marge de gauche en regard des « suites exactes
$A_n^p \to A_n^q \to M_n \to 0$ » : « cf. EGA IV 18.3.2.1. » (la référence
même que cite sa réponse, page 97) ; en outre un double trait vertical, de
main non établie, dans la marge de gauche en regard du dernier alinéa}

\page{96}

\note{lettre dactylographiée de lui, datée à la main en haut à droite. Les
flèches, les signes $\in$, $\leq$ et le tilde de $\widetilde{M}$ ont été
laissés en blanc par la machine et jamais complétés ; on les rétablit dans
les formules là où le texte les fixe. Les surfrappes ne sont pas relevées}

15.9.67.

Cher Dieudonné,

Tes objections de ta lettre du 11 Septembre sont encore fondées. D'ailleurs,
il ne me semble pas évident que le fait de pouvoir trouver $p,q$ fixes tels
qu'on ait des suites exactes $A_n^p \to A_n^q \to M_n \to 0$, permette de
s'en tirer ; tant mieux si tu y arrives. L'existence de ces $p$ et $q$ me
semble d'autre part facile, en utlisant le

\emph{Lemme} \quad Soit $A$ un anneau, $M$ un $A$-module,
$(f_i)_{1 \leq i \leq n}$ des éléments de $A$ engendrant l'idéal unité,
supposons que pour tout $i$, le $A_{f_i}$-module $M_{f_i}$ soit engendré par
$m_i$ générateurs. Alors $M$ est engendré par $m \quad m_i$
générateurs.\note{entre $m$ et $m_i$, un blanc que la machine n'a pas rempli
et que la page ne permet pas de fixer ; la suite (« $mp$ éléments, où $m$ est
le nombre des $f_i$ ») suggère $\sum m_i$}

Démonstration : soient $g^i_j \in M_{f_i}$ $(1 \leq j \leq m_i)$ les
générateurs de $M_{f_i}$, on aura donc $g^i_j = h^i_j / f_i^{N_{ij}}$, avec
$h^i_j \in M$. Alors les $h^i_j$ pour $j$ fixes engendrent $M$ sur l'ouvert
$\operatorname{Spec}(A_{f_i})$ de $\operatorname{Spec}(A_f)$, donc comme pour
$i$ varaible ces ouverts recouvrent $\operatorname{Spec}(A)$, il s'ensuit que
les $h^i_j$ pour $i,j$ variables engendrent $M$, donc $M$, ce qui établit le
lemme.

Revenant alors à la situation de 10.10.5, on sait que $M = \lim M_n$ est un
$A$-module de type fini, considérons alors un épimorphisme
$u : A^q \to M$. Soient $f_i \in A$ dont les images dans $A_0$ engendrent
l'idéal unité, et tel que sur les ouverts correspondants $X_i$ de
$X = \operatorname{Specf}(A)$, $F$ admette une présentation finie, donc soit
de la forme $\widetilde{M_i}$, où $M_i$ est un $A_i = A_{f_i}$-modul\supplied{e} de
présentation finie. Alors $u$ restreint à $X_i$ définit un épimorphisme
$A_i^q \to M_i$, et comme $M_i$ est de présentation finie, il existe un
homomorphi\supplied{s}me $v_i$ rendant exacte la suite $A_i^p \to A_i^q \to M_i \to 0$.
Tensorisant par $A_0$,
on en déduit une suite exacte
$(A_n)_{f_i}^p \to (A_n)_{f_i}^q \to (M_n)_{f_i} \to 0$, ce qui montre que,
si $R_n = \operatorname{Ker}(A_n^q \to M_n)$, alors $R_n$ est sur
$\operatorname{Spec}(A_{n\, f_i})$ engendré par $p$ éléments. Donc en vertu
du lemme $R_n$ est engendré par $mp$ éléments
\note{« $A_i =$ » est tapé dans l'interligne au-dessus de « $A_{f_i}$ » ;
« Tensorisant par $A_0$ » est ainsi à la machine, la suite porte sur $A_n$}

\page{97}

\note{suite de la même lettre, second feuillet ; pas de signature}

où $m$ est le nombre des $f_i$, et $mp$ est bien indépendant de $n$.

La difficulté qui semble rester est de trouver les suites exactes que tu
demandes de telle façon qu'elles se recollent, pour $n$ variable. Donc, $u$
étant déjà choisi, de trouver un homomorphisme $A^p \to A^q$ dont l'image soit
$\operatorname{Ker} u$ \ldots\ Si tu y arrives, on pourrait présenter encore
10.10.5 sous forme de trois conditions équivalentes, mais en demandant dans
b) une « prsentation finie uniformément en $n$ ». Sinon, il faudrait trouver
un contre-exemple à l'implication a) $\Rightarrow$ c), car il me semble que
la question devrait être tirée au clair. En tout état de cause, il faudra
donner, en corollaire, l'équivalence des condition\supplied{s} suivantes, valables sans
hypothèses noethériennes :

\begin{itemize}
\item[a)] $F$ est localement libre de type fini.
\item[b)] $F$ est isomorphe à la limite projective d'une suite $(F_n)$ de
$\underline{O}_{X_n}$-Modules loc libres de type fini qui se recollent.
\item[c)] $F$ est isomorphe à un $\widetilde{M}$, avec $M$ un $A$-module
projectif de type fini.
\end{itemize}

Pour la démonstration, on procède comme dans 10.10.5 en utilisant EGA IV
18.3.2.1. Ce lemme devrait d'ailleurs venir en corollaire après
$0_{\mathrm{I}}$ 7.2.9.

Pour les modifications que je préconaisais pour I 10.11, elles tombent à
l'eau si on n'arrive pas à arranger 10.10.5 sans hypothèses noethériennes ;
on peut cependant dire que si $F$ sur $X$ est de présentation finie, alors il
est limite projective de faisceaux $F_n$ de présentation finie sur les $X_n$
qui se recollent (mais on n'a pas une réciproque), et que $F$ est localement
libre \uncertain{ssi} les $F_n$ le sont. Et les autres énoncés du n°
10.11. restent valables sans hypothèse noethérienne, sauf 10.7.11.2 et la
partie « injectivité » dans 10.11.9 (sauf erreur).
\note{« ssi » : frappe surchargée, qui se lit « sis » ou « sàs »}

\section*{EGA II}

\note{inscrit au crayon en haut à droite du feuillet de garde (page 98), qui
ne porte rien d'autre ; vraisemblablement de sa main. Le feuillet ne reçoit
pas de numéro de page}

\page{99}

\note{à l'encre, sur une feuille à en-tête de l'IHÉS ; les lettres a) à d)
sont entourées, les numéros des volumes d'EGA soulignés}

\begin{itemize}
\item[(a)] transférer EGA IV 4.8

\uncertain{Également} \uncertain{écrit} des morphismes affines de
\uncertain{préschémas} formels.\note{ces trois lignes sont tenues par une
accolade à gauche}

Appendice \uncertain{pour} \ill{} aux schémas formels !\note{« pour » est
récrit d'un trait épais ; un petit trait vertical en marge gauche de la
ligne}

\item[(b)] Cf EGA II 6.5.10 \quad ajouter formules dans le \ill{}

6.5.11 \quad ajouter formule pour le cas où $X' \to X$ est gén.\
\uncertain{décomposé} (\ill{} \uncertain{moins} \ill{} !)

6.5.7 \quad renvoi à EGA IV 21.

\item[(c)] Remonter EGA IV 19.4.11 (revu et augmenté)

\item[(d)] 7.4.\uncertain{2}. \quad \uncertain{Admettre} \ill{} vide ---
\end{itemize}

\end{document}
